\documentclass[../main.tex]{subfiles}
\begin{document}
%==========================================TIÊU ĐỀ TẬP========================================
\begin{table}[h]
    \begin{adjustwidth}{-1cm}{}
        \begin{tabular}{>{\centering\arraybackslash}p{6cm}|>{\raggedright\arraybackslash}p{12.5cm}}
            \multirow{5}{6cm}
            % Cột bên trái: ÉPISODE và số tập
            {\\[-40pt]
                \flaregothic\fontsize{20pt}{20pt}\selectfont \centering
                \color{eptype}HISTOIRE\color{black}\\
                \freshbot\fontsize{64pt}{64pt}\selectfont
                \color{epnum}E\\[6pt]
                \cafeteria\fontsize{20pt}{20pt}\selectfont\color{epnum}(D-E)\color{black}
                \\[18pt]
                % \sen\fontsize{18pt}{18pt}\selectfont
                % \color{epattr}Interstitiel
                \color{black}
            }
             &
            % Dòng thứ nhất: tên tập tiếng Nhật
            {\fotskip\fontsize{18pt}{18pt}\selectfont \ruby{三}{さん}\ruby{階}{かい}の\ruby{小}{ちい}さな\ruby{城}{しろ}
            }                                                      \\[6pt]
            % Dòng thứ hai: tên tập tiếng Anh
             & {\cafeteria\fontsize{18pt}{18pt}\selectfont A Castle on the Third Floor}                                                                       \\[4pt]
            % Dòng thứ ba: tên tập tiếng Việt
             & {\vnmsans\fontsize{18pt}{18pt}\selectfont Tòa lâu đài tầng ba}                                                                                 \\[8pt]
            % Dòng thứ tư: Nhân vật xuất hiện
             & {\caviar{\fontsize{14pt}{14pt}\selectfont Nhân vật: \emp{kuon2} \emp{ryuuhou2} \emp{kaiho2} \emp{suzuno2} \emp{naomi2} \emp{coco} \emp{aloe}}}
            \\[8pt]
        \end{tabular}
    \end{adjustwidth}
\end{table}
%=========================================TRUYỆN CHÍNH========================================
\color{black}

\txtbx[2]{
    {\color{vol} \uline{Tóm tắt:}} Sau những đêm đầu tiên đại gia đình chen chúc ngủ chung ấm cúng trên sàn phòng khách tầng năm, bố Ryujin quyết định chấn chỉnh nếp nhà bằng cách cải tạo lại toàn bộ tầng ba—vốn là tầng để trống từng cho thuê và chỉ dùng làm nơi tổ chức tiệc—thành không gian sinh hoạt riêng biệt cho phái nữ. Không ồn ào gạch vữa hay máy móc tự động kỳ dị, cả gia đình cùng Coco và Aloe xắn tay áo bài trí, lắp ráp giường tầng, bàn học đôi và từng tấm vách ngăn gỗ ấm áp. Nhưng khi căn phòng hoàn thiện khang trang, cô bé út Naomi lại ôm chặt chú chó già ngấn lệ hỏi: ``Vậy là\dots\ Naomi không được ngủ chung với Onii-chan nữa sao ạ?''. Bằng sự dịu dàng của người anh cả và món quà kết nối bất ngờ, khoảng cách giữa các tầng lầu bỗng chốc hóa thành sợi dây gắn kết thiêng liêng.
}

\pov{Hikawa Kuon}

Sau ngày làm thủ tục xét nghiệm ADN và nhận giấy tờ hộ tịch đầy nhiệm màu ở trụ sở ủy ban phường Anwa, căn nhà của chúng tôi chính thức bước vào một trang hoàn toàn mới. 

Thế nhưng, niềm vui sum họp quá đỗi đột ngột ấy cũng kéo theo những tình huống dở khóc dở cười mà trước đây tôi chưa từng nghĩ tới.

Cần phải nói rõ rằng, căn phòng sinh hoạt chung kiêm chỗ ngủ của bố con tôi ở nửa bên phải tầng năm vốn là một không gian hình chữ L với diện tích khá khiêm tốn: vỏn vẹn chưa đầy năm mươi mét vuông. Trước đây, nơi này chỉ có hai bố con tôi ở thì tương đối vừa vặn: một góc kê chiếc tivi ba mươi lăm inch cũ kỹ, một bộ bàn trà gỗ, tủ kính đựng đồ lưu niệm, chiếc quạt cây hỏng nút bấm phải chỉnh bằng tuốc-nơ-vít, và góc bàn làm việc ngổn ngang dây cáp, màn hình máy tính của tôi. 

Vậy mà suốt ba đêm liền, theo lệnh của bố tôi, cả bảy con người lớn bé—gồm bố mẹ, tôi, Ryuuhou, cùng ba cô em gái mới đón về là Kaiho, Suzuno, Naomi—cộng thêm cả ba mẹ con nhà chó là con Mi già nua hơn mười ba tuổi, chú chó Tô hiếu động và chú cún con Ki lí lắc, đều dồn hết vào đây, trải đệm mút nằm la liệt khắp sàn nhà hình chữ L!

Ý định của bố thì rất đẹp: để các con gắn kết tình thân sau bao năm trời thất lạc. Nhưng trên thực tế, việc nhồi nhét ngần ấy sinh mạng vào năm mươi mét vuông đã biến căn phòng thành một ``bãi chiến trường'' thu nhỏ.

Ban đêm, những tấm đệm trải kín mít không còn lấy một khe hở để đặt chân. Naomi ngủ say thì ôm chặt lấy cánh tay tôi như một chú gấu túi nhỏ, hễ tôi cựa quậy rút tay ra là bé con lại giật mình dụi mắt mếu máo; chú chó Mi già cuộn tròn ngay dưới bụng con bé sưởi ấm, trong khi hai con cún Tô và Ki thì thi nhau rúc vào chăn liếm mặt Ryuuhou. Khổ sở nhất là Ryuuhou: con bé ngủ có tật xấu quay mòng mòng như compa, nửa đêm đạp tung cả chăn rồi vắt thẳng hai cẳng chân ngang qua sườn tôi. Suzuno nằm nép cạnh mẹ Kaien sát vách tường, nửa đêm muốn dậy đi uống nước mà cứ loay hoay ngồi im thin thít, chẳng dám nhúc nhích vì sợ giẫm phải người khác hay đạp trúng đuôi chó.

Đến sáng ngày thứ tư thì mọi sự bất tiện chính thức bùng nổ thành một cơn hỗn loạn đỉnh điểm.

*RẦM! RẦM! RẦM!*

Tiếng đập cửa dồn dập vang lên inh ỏi từ phía góc nhà vệ sinh duy nhất của tầng năm. Kaiho trong bộ trang phục chỉnh tề, mái tóc dài buộc gọn sau gáy, hai tay chống nạnh, mặt mày đanh lại gõ mạnh vào cánh cửa gỗ:

``Hikawa Ryuuhou! Em ở trong đó bôi trét cái gì mà hơn hai mươi phút rồi chưa chịu ra thế hả?! Chị sắp trễ giờ nộp báo cáo đồ án rồi đây này!''

Bên trong, tiếng nước xối rột rẹt hòa cùng tiếng hét the thé vọng ra không hề kém cạnh:

``Em đang chải lại cái đầu tổ quạ này chứ có bôi trét gì đâu! Chị Kaiho chờ năm phút nữa đi, ai bảo chị không dậy sớm hơn người ta!''

``Em nói cái gì?! Chị dậy từ sáu giờ sáng, nhưng vừa bước ra khỏi đệm đã bị chân con Mi ngáng ngã, lại còn phải đợi mẹ rửa mặt xong nữa! Mở cửa ra ngay!''

Tôi đứng cạnh đó, tay cầm chiếc bàn chải đánh răng sùi bọt mép, một bên mắt còn díp lại vì thiếu ngủ, chỉ biết lắc đầu thở dài ngao ngán: ``Hai đứa bớt ồn ào giùm anh một tí đi\dots\ Anh còn phải chuẩn bị tài liệu để lát nữa bắt chuyến tàu sớm sang văn phòng Tokyo nữa này\dots''

Suzuno đứng nép một bên hành lang, hai tay ôm chiếc khăn mặt nhỏ, giọng em lí nhí đầy bẽn lẽn: ``Anh Kuon\dots\ chị Kaiho\dots\ hay là để em xuống tầng bốn nhờ nhà vệ sinh bên phòng bà nội với bác Hijaki ạ?''

``Không được đâu Suzu!'' Kaiho lập tức quay sang gạt đi. ``Bên tầng bốn chia đôi, phòng khách bác Hijaki một bên, phòng bà nội một bên. Bác với bà đều có tuổi rồi, sáng sớm kéo nhau xuống đó ầm ĩ thì còn ra thể thống gì nữa!''

Đúng lúc đó, từ góc đệm ngủ phía trong, bé út Naomi trong bộ đồ ngủ in hình gấu bông lững thững bước ra, hai tay ôm chú cún con Ki lông trắng muốt. Đôi mắt xanh ngọc tròn xoe của cô bé ngơ ngác chớp chớp, rồi chỉ tay xuống sàn nhà, cất giọng ngây thơ:

``Chị Kaiho ơi\dots\ Ki vừa gặm mất một chiếc nơ trên giày da của chị rồi ạ\dots''

Kaiho quay phắt lại. Nhìn chú cún con Ki mới được vài tháng tuổi đang ngoe nguẩy đuôi nhay rách mép chiếc nơ giày công sở của mình, cô nàng ôm trán, ngửa mặt lên trần nhà kêu trời:

``Trời đất ơi là trời! Năm mươi mét vuông mà nhét bảy con người với ba con chó, cái nhà này sắp biến thành cái rạp xiếc trung ương rồi!''

Giữa lúc tiếng cãi cọ, tiếng cún con ẳng ẳng và tiếng thở dài hòa vào nhau thành một mớ âm thanh hỗn độn, bố Ryujin từ ngoài ban công bước vào. Chứng kiến toàn bộ cảnh tượng nháo nhào ấy, ông bưng chén trà đặc trên tay, đặt mạnh xuống mặt bàn gỗ đánh một tiếng *CỘP* thật giòn.

Cả căn phòng lập tức nín bặt.

Bố tôi vuốt chòm râu quai nón, đảo mắt nhìn khắp lượt bảy con người đang đứng chen chúc giữa căn phòng chật chội, rồi cất giọng trầm vang đầy uy lực của người chủ gia đình:

``Thôi, đủ rồi! Không thể để cái cảnh này tiếp diễn thêm một ngày nào nữa! Nhà Hikawa chúng ta sống ở đất Kawachi này bao đời nay có nếp nhà gia phong đàng hoàng. Các con gái đều đã lớn tướng cả rồi, nam nữ thanh thiếu niên ai lại chen chúc chung chạ trong một cái phòng chưa đầy năm mươi mét vuông thế này! Vừa mất phép tắc, vừa nhếch nhác loạn cào cào cả lên! Hôm nay là ngày nghỉ, bố lệnh cho cả nhà tổng động viên: dọn dẹp, cải tạo lại toàn bộ tầng ba đang để trống làm phòng riêng cho hội phụ nữ!''

\ct

Tầng ba của ngôi nhà chúng tôi vốn là một khoảng không gian khá đặc biệt.

Trước đây, khi căn nhà sáu tầng này mới hoàn thành, tầng ba từng được bố mẹ cho một văn phòng đại diện nhỏ thuê lại để kiếm thêm thu nhập. Sau khi họ hết hợp đồng chuyển đi, nơi này được để trống hoàn toàn. Ngoại trừ chiếc bàn gỗ dài và hàng ghế tựa đặt ở giữa phòng dùng để tụ họp đại gia đình ăn đại tiệc vào những dịp giỗ chạp hay lễ tết, phần lớn diện tích xung quanh chỉ dùng để chứa vài thùng các-tông và đồ đạc ít dùng.

Lần này, chúng tôi không hề có ý định nhờ cậy đến cỗ máy xây dựng tự động kỳ dị của bác Hijaki - thứ máy móc từng dựng nên căn phòng nghiêng bốn mươi lăm độ dở khóc dở cười cho Coco dạo trước. Công việc hôm nay thuần túy là cải tạo nội thất: dọn dẹp bụi bặm, phân chia công năng, và biến một mặt bằng lạnh lẽo thành một chốn nghỉ ngơi ấm cúng, riêng tư nhất cho mẹ Kaien và bốn cô em gái.

``Nào, phân công nhiệm vụ nhé!'' Tôi đứng giữa sảnh tầng ba, hai tay cầm hai cuộn thước dây, bắt đầu vạch kế hoạch. ``Bố và bác Hijaki sẽ phụ trách kiểm tra lại đường điện và hệ thống thông gió. Em và Kaiho sẽ lắp ráp các khung vách ngăn gỗ thông để chia không gian ngủ. Mẹ, Ryuuhou và Suzuno phụ trách lau dọn bụi bặm, còn Naomi\dots''

Tôi cúi xuống nhìn cô em út đang lon ton ôm theo chiếc chổi lông gà nhỏ xíu: ``Nhiệm vụ tối cao của Naomi là làm đội trưởng giám sát công trình, cùng Mi-chan kiểm tra xem các anh chị có làm việc chăm chỉ không, được chứ?''

Naomi nghe vậy liền đứng nghiêm, giơ bàn tay nhỏ xíu lên trán làm động tác chào kiểu hải quân: ``Rõ ạ! Naomi sẽ canh chừng thật cẩn thận!''

Tiếng cười rộn rã lập tức xua tan đi sự lạnh lẽo lâu ngày của tầng ba.

Đúng lúc đó, cánh cửa cầu thang bật mở, một giọng nói sang sảng đầy năng lượng cất lên:

``Alo alo! Có ai cần một con rồng vận chuyển hạng nặng không nào?''

Coco bước vào, mặc chiếc áo phông ngắn tay khỏe khoắn, theo sau là Aloe với chiếc tạp dề hoa nhã nhặn xách theo làn đồ ăn nhẹ.

``Nghe nói nhà Hikawa hôm nay có chiến dịch đại tu phòng ốc, sao thiếu hai đứa này được!'' Coco cười lớn, vỗ bồm bộp vào bắp tay săn chắc. ``Cứ có món đồ nào nặng trên một tạ thì cứ giao hết cho rồng thần này nhé!''

``May quá, đúng lúc bọn anh đang cần chuyển chiếc tủ quần áo gỗ ba cánh từ tầng năm xuống,'' tôi cười đáp.

Và quả thực, sức khỏe của một nàng rồng không phải trò đùa. Trong khi tôi và bố phải hò hét lấy thế mới nhích nổi chiếc tủ gỗ gụ nặng trịch, thì Coco chỉ cần một tay nâng đáy, một tay giữ mép tủ, thong thả bước từng bước xuống cầu thang gỗ nhẹ tênh như đang vác một thùng bìa các-tông rỗng. Chỉ có điều, chiếc tủ hơi cồng kềnh khiến nó suýt quẹt vào lớp sơn tường hành lang, làm Kaiho đứng bên cạnh phải vội vàng giơ tay đỡ mép gỗ, chau mày nhắc nhở:

``Này nàng rồng ngổ ngáo, đi chậm thôi! Quẹt xước góc tường là Aniki bắt cô đền tiền sơn đấy!''

Coco lè lưỡi tinh nghịch: ``Biết rồi, biết rồi mừ Kaiho-paisen! Em có đo đạc khoảng cách bằng mắt rồng rồi, chuẩn từng mi-li-mét luôn!''

Trong khi đó, ở góc đối diện, Aloe đang cùng mẹ Kaien và Suzuno tỉ mỉ đo từng khung cửa sổ nhìn ra dải công viên ven sông. Nàng succubus khéo léo dùng kéo bấm cắt từng đoạn rèm cửa màu kem sữa mềm mại, rồi luồn vào thanh treo bằng đồng sáng bóng.

``Rèm màu này hợp lắm, Kaien-san,'' Aloe mỉm cười dịu dàng, đôi tai nhọn khẽ rung rinh theo nhịp tay đưa kim. ``Vừa cản được ánh nắng gay gắt buổi trưa, vừa giữ được sự ấm áp và thoáng đãng khi gió sông thổi vào buổi chiều.''

Suzuno đứng bên cạnh phụ luồn móc treo, đôi mắt hai màu lấp lánh sự ngưỡng mộ: ``Chị Aloe khéo tay thật đấy\dots\ Sau này chị chỉ cho em cách may viền ren này với nhé?''

``Ừm, bất cứ lúc nào Suzuno-chan thích!'' Aloe vui vẻ gật đầu.

Ở khu vực trung tâm, tôi và Ryuuhou cặm cụi lắp ráp hai chiếc giường tầng bằng gỗ thông tự nhiên. Tiếng máy bắt vít cầm tay *rè\dots\ xạch* vang lên giòn giã. Tôi cẩn thận dùng giấy ráp chà mịn từng góc cạnh của thanh giằng gỗ, rồi dán thêm các miếng mút cao su mềm vào các góc nhọn tiếp giáp.

Ryuuhou vừa phụ tôi giữ thanh giằng vừa lau mồ hôi trên trán, tò mò hỏi: ``Onii, dán mấy miếng mút này làm gì cho mất công thế? Giường gỗ để mộc nhìn mới ngầu chứ.''

Tôi gõ nhẹ sống tay vào trán cô em: ``Ngầu cái gì mà ngầu. Tối Naomi hay ngái ngủ lơ mơ đi tìm dép, không bọc lại lỡ con bé cộc đầu vào góc giường thì sưng vù trán à?''

Ryuuhou sờ trán mình, bĩu môi nhưng ánh mắt lại lấp lánh vẻ phục lăn: ``Đúng là anh cả\dots\ chu đáo đến phát sợ luôn.''

Không gian tầng ba dần dần biến đổi kỳ diệu dưới bàn tay của cả gia đình. Chiếc bàn gỗ dài trước đây vốn dùng cho các bữa tiệc lớn nay được chúng tôi lau chùi bóng loáng, kê sát về phía góc ban công ngập tràn ánh sáng. Kuon tôi đặt thêm một tấm khăn trải bàn kẻ ca-rô trang nhã, biến nơi này thành khu vực phòng khách nhỏ dành riêng cho năm mẹ con thưởng trà, đọc sách hay trò chuyện mỗi tối mà không hề làm mất đi công năng tổ chức tiệc gia đình khi có dịp lễ hội lớn.

Ngay bên cạnh, tôi và Kaiho thiết kế một chiếc bàn học đôi bằng gỗ sồi sáng màu, có giá sách hai tầng phía trên cho Ryuuhou và Suzuno cùng ngồi học bài. Giữa căn phòng, bức vách ngăn bằng nan gỗ thông thoáng khí được dựng lên khéo léo, chia tách thành hai khu vực ngủ riêng biệt: một góc yên tĩnh bên trong dành cho mẹ Kaien, và một gian phòng ngủ rộng rãi, thoáng mát phía ngoài với hai chiếc giường tầng chắc chắn dành cho bốn chị em gái.

Dưới chân giường tầng dưới, một chiếc đệm êm ái nhỏ xíu có in hình bàn chân cún con cũng được đặt ngay ngắn—chỗ ngủ dành riêng cho chú chó Mi-chan của Naomi.

\ct

Hoàng hôn dần buông xuống khu phố Anwa. Những tia nắng vàng cam cuối ngày xuyên qua lớp rèm voan màu kem sữa, rải một lớp ánh sáng lộng lẫy và ấm áp lên toàn bộ không gian tầng ba.

Căn phòng thơm ngát mùi gỗ thông mới quyện với hương tinh dầu sả chanh dịu nhẹ mà Aloe vừa xông. Đèn chùm nhỏ với ánh sáng vàng ấm trên trần được bật lên, khiến cả tầng ba trông chẳng khác nào một góc của căn hộ cao cấp tràn đầy sức sống.

``Xong rồi! Hoàn hảo ngoài sức tưởng tượng!'' Ryuuhou thả phịch người xuống chiếc ghế bành ở góc bàn trà, dang hai tay thở phào nhẹ nhõm.

Bố tôi vừa nhấp ngụm nước chè tươi vừa gật gù khen ngợi: ``Khá lắm! Vừa ngăn nắp, vừa ấm cúng mà lại giữ được nếp nhà gia phong. Từ tối nay, hội con gái chính thức chuyển xuống đây sinh hoạt!''

Mẹ Kaien, Kaiho và Suzuno đều vô cùng hào hứng bắt đầu mở các thùng các-tông để xếp quần áo, sách vở và đồ đạc cá nhân vào các ngăn tủ mới. Tiếng cười đùa râm ran khắp căn phòng.

Thế nhưng, giữa bầu không khí náo nhiệt ấy, tôi chợt nhận ra có một điều gì đó bất thường.

Tôi quay đầu nhìn quanh, và bắt gặp một bóng dáng nhỏ bé đang đứng nép sát vào mép cánh cửa ra vào.

Là Naomi.

Cô bé ôm ghì lấy chú chó Mi vào lòng, đôi bàn tay nhỏ bé siết chặt lớp lông xơ xác của con chó già. Đôi mắt xanh lục bảo to tròn vốn luôn trong veo và rạng rỡ, lúc này bỗng chốc ầng ậng nước. Đôi môi nhỏ mím chặt lại, bờ vai gầy khẽ run rẩy từng nhịp.

Nụ cười trên môi tôi lập tức chững lại. Tôi đặt chiếc tuốc-nơ-vít xuống bàn, chậm rãi bước lại gần cô em út, rồi quỳ một gối xuống sàn để tầm mắt ngang bằng với em: ``Naomi à? Sao thế em? Căn phòng mới có chỗ nào làm em không thích sao?''

Nghe tiếng tôi hỏi, giọt nước mắt kìm nén nãy giờ của cô bé rốt cuộc lăn dài trên má. Naomi lắc đầu nguầy nguậy, rồi cúi gằm mặt xuống lớp lông của Mi-chan, giọng em nghẹn ngào, run rẩy và nhỏ như tiếng muỗi kêu:

``Không phải đâu ạ\dots\ Phòng\dots\ phòng đẹp lắm, giống như lâu đài của công chúa vậy\dots''

``Thế sao công chúa của anh lại khóc?'' Tôi dịu dàng đưa ngón tay gạt đi giọt nước mắt ấm nóng trên gò má em.

Naomi ngước lên, đôi mắt đỏ hoe nhìn tôi đầy sợ hãi và buồn bã: ``Nhưng mà\dots\ bố bảo từ tối nay con gái phải ngủ ở tầng ba\dots\ Vậy là\dots\ Naomi không được ngủ chung với Onii-chan nữa sao ạ?''

Lời nói non nớt cất lên khiến cả căn phòng tầng ba bỗng chốc im bặt.

Tiếng gấp quần áo của mẹ Kaien dừng lại. Kaiho đang cầm móc áo cũng đứng sững người. Bố tôi ngơ ngác gãi đầu, trong khi Ryuuhou và Suzuno đều đồng loạt quay lại nhìn cô em út.

Trong đôi mắt ngây thơ của đứa trẻ vừa trải qua muôn vàn biến cố ấy, tôi đọc được tất cả nỗi sợ hãi mà em đang chôn giấu.

Suốt nhiều năm trời lưu lạc trong không gian song song đầy bóng tối và sự cô độc, Naomi chỉ có một mình con chó già làm bạn. Khi được đưa về căn nhà này, được nằm ngủ trên tầng năm bên cạnh tôi và bố mẹ, cảm nhận được hơi thở ấm áp, tiếng ngáy khẽ của bố và tấm lưng vững chãi của người anh cả ngay bên cạnh, cô bé mới lần đầu tiên cảm nhận được thế nào là sự an toàn tuyệt đối. Em sợ rằng nếu phải rời khỏi tầng năm, nếu phải tách xa người anh trai mà em xem như cột trụ che chở, thì cái hơi ấm kỳ diệu ấy sẽ vụt biến mất như một giấc mơ thoáng qua.

Cổ họng tôi dâng lên một cảm giác xót xa khôn tả.

Tôi khẽ mỉm cười, đưa tay xoa nhẹ mái tóc mềm của Naomi, rồi kéo cả cô bé lẫn chú chó Mi vào một vòng ôm thật ấm áp: ``Ngốc ạ. Ai bảo với Naomi là xuống tầng ba thì không còn được ở gần Onii-chan nữa nào?''

Naomi thút thít, tựa cằm vào vai tôi: ``Nhưng\dots\ Onii-chan ngủ ở tầng năm cơ mà\dots\ Tầng ba xa lắm\dots''

``Xa cái gì mà xa,'' tôi bật cười, buông em ra rồi thò tay vào túi áo khoác, lấy ra một chiếc hộp nhỏ bằng nhựa mica mà tôi đã cặm cụi hàn mạch điện trưa nay.

Đó là một chiếc đèn ngủ mini hình ngôi sao năm cánh phát ra ánh sáng vàng dịu, bên dưới có gắn một chiếc loa nhỏ và một chiếc nút bấm màu xanh ngọc.

Tôi đặt chiếc đèn ngôi sao vào hai bàn tay nhỏ bé của Naomi, ân cần giải thích: ``Em nhìn xem, chiếc đèn ngôi sao này anh đã nối một đường dây tín hiệu riêng chạy thẳng lên bàn làm việc của Onii-chan ở tầng năm đấy. Mỗi tối trước khi đi ngủ, Naomi chỉ cần bật đèn lên là trên phòng anh cũng sẽ sáng một ngôi sao nhỏ tương tự.''

Naomi chớp chớp mắt, ngơ ngác nhìn món đồ chơi kỳ lạ: ``Thật ạ?''

``Thật một trăm phần trăm!'' Tôi gõ nhẹ vào chóp mũi em. ``Và chiếc nút này là bộ đàm khẩn cấp. Bất cứ khi nào Naomi thấy sợ bóng tối, nhớ anh, hay chỉ đơn giản là muốn nghe kể chuyện cổ tích, em chỉ cần bấm giữ nút này và gọi: `Onii-chan ơi!'. Anh ở ngay trên trần nhà em thôi, chỉ cần chạy ba bước xuống cầu thang là có mặt bên cạnh em ngay lập tức. Được chưa nào?''

Naomi nhìn chiếc đèn ngôi sao phát sáng lấp lánh trong tay, rồi ngước nhìn trần nhà tầng ba. Nỗi sợ hãi vô hình trong mắt cô bé dần tan biến, thay vào đó là một tia sáng kỳ diệu.

Kaiho bước tới từ phía sau, khoanh hai tay trước ngực, giả vờ hất cằm nhưng ánh mắt lại ngập tràn sự cưng chiều:

``Hơn nữa, ở căn phòng này Naomi có cả mẹ Kaien, có chị Kaiho, chị Suzuno và chị Ryuuhou ở bên cạnh cơ mà. Bốn người bọn chị bảo vệ em còn không đủ hay sao? Tên Aniki kia mà dám chậm trễ không xuống kể chuyện cho em, chị sẽ đích thân phi dép lên tầng năm lôi cổ anh ta xuống!''

Ryuuhou cũng nhảy bổ lại, bẹo nhẹ hai má bánh bao của Naomi: ``Đúng đấy! Ở với ba anh em trai bọn anh chán ngắt, xuống đây tha hồ chơi búp bê với nghe nhạc rock cùng chị Ryuu có sướng hơn không!''

Suzuno mỉm cười dịu dàng, nhẹ nhàng nắm lấy bàn tay còn lại của cô bé: ``Chị sẽ đọc truyện tranh cho Naomi nghe trước khi ngủ nhé?''

Đến lúc này, Naomi mới quệt sạch những giọt nước mắt còn vương trên khóe mi. Nụ cười rạng rỡ như ánh nắng ban mai lại nở bừng trên gương mặt thiên thần nhỏ. Cô bé ôm chặt chiếc đèn ngôi sao vào ngực, cất giọng trong trẻo:

``Vâng ạ! Naomi thích căn phòng tầng ba lắm! Cảm ơn Onii-chan, cảm ơn mọi người nhiều ạ!''

Bố Ryujin đứng nhìn cảnh tượng ấy, thở phào một hơi dài rồi cười ha hả: ``Thế là êm xuôi nhé! Tối nay bố thưởng cho cả nhà một nồi lẩu bò hoành tráng ăn mừng tân gia tầng ba luôn!''

Cả căn phòng lại vang dội tiếng vỗ tay và reo hò phấn khích.

\ct

Đêm hôm ấy.

Sau bữa lẩu bò no nê và những trận cười nghiêng ngả, cả khu phố Anwa chìm dần vào giấc ngủ tĩnh mịch.

Phòng khách tầng năm bỗng chốc trở nên thênh thang lạ thường. Không còn cảnh những tấm nệm trải dài la liệt, không còn những chiếc chân đạp lung tung hay tiếng thở ríu rít của đàn em. Chỉ còn lại hai bố con tôi nằm trên hai chiếc giường đơn quen thuộc. Không gian yên ắng đến mức có thể nghe rõ cả tiếng gió đêm thổi qua rặng cây bên công viên bờ sông.

Tôi nằm nghiêng người trên giường, nhìn chiếc màn hình máy tính vẫn còn sáng mờ mờ trên bàn làm việc.

Đúng lúc ấy, một tiếng *tút\dots\ tách* thật khẽ vang lên từ chiếc máy bộ đàm mini đặt cạnh bàn phím.

Ánh đèn ngôi sao trên bàn khẽ nhấp nháy một màu xanh ngọc dịu dàng.

Rồi một giọng nói lí nhí, ngọng nghịu nhưng ngọt ngào như kẹo sữa truyền qua loa:

``Onii-chan ơi\dots\ Naomi với Mi-chan chúc Onii-chan ngủ ngon ạ\dots''

Tôi bất giác mỉm cười, cảm thấy một luồng hơi ấm ngập tràn lồng ngực. Tôi vươn tay bấm nút đáp lại:

``Ừ, Naomi ngoan. Chúc em và Mi-chan ngủ ngon nhé.''

Tiếng lách cách ngắt kết nối vang lên, rồi căn phòng lại trở về với sự tĩnh lặng bình yên.

Tôi nhắm mắt lại, lòng thanh thản đến lạ kỳ.

Tầng ba giờ đây đã có mẹ và các em gái, tầng năm có tôi và bố, tầng bốn có bà và bác Hijaki, tầng hai có Aloe, tầng bảy có Coco. Mỗi tầng lầu là một góc trời riêng, nhưng tất cả đều chung một nhịp đập dưới mái nhà này. Khoảng cách giữa các tầng lầu chẳng những không làm chúng tôi xa nhau, mà trái lại, nó càng chứng minh rằng ngôi nhà Hikawa đã thực sự trở thành một bến đỗ vững chãi, ấm áp và trọn vẹn nhất cho tất cả chúng tôi.

\end{document}
